\section{Introduction}\label{sec:intro}

Reconstruction on trees asks whether the spin at the root can be inferred
with a nonvanishing advantage from the spins on arbitrarily distant
levels. The Kesten--Stigum second-eigenvalue scale originates in
multitype branching-process theory \cite{KS66}; broadcasting and
information flow on trees were developed in
\cite{EKPS00,MP03}. For the binary Ising model, nonlinear potential
and capacity furnish finer criteria on nonregular trees, including at
criticality. The general-tree capacity theory of Pemantle and Peres
\cite{PP10} is an important antecedent: in its Ising setting, an
$L^3$ capacity occurs for plus boundary conditions. Robust
reconstruction and the role of the second eigenvalue were investigated
by Janson and Mossel \cite{JM04}.

For the symmetric Potts broadcast channel, the transition law is a
convex combination of the identity and the uniform kernel. The
four-state case lies at a cancellation point in the posterior
expansion; the standard quadratic correction is insufficient and a
higher-order term controls the critical regime. The sharpness and
nonsharpness phenomena across the number of states were studied by
Sly \cite{Sly11}, with later results of Mossel, Sly and Sohn
\cite{MSS25}. Two recent manuscripts \cite{OAIThreshold,OAICapacity}
provide, respectively, a rigorously certified four-state posterior
moment inequality and an $L^3$ criterion for bounded-degree,
constant-channel deterministic trees. The present note extends the
latter to heterogeneous channels and removes the bounded-degree
restriction, while preserving its explicit dependence on the former.

Our two principal transfers are different. The 
\emph{necessity} argument modifies the information variable by
$f(m)=m+m^3/3000$. The moment-saving inequality makes this variable
subadditive under every binary combination, so sequential products
can be controlled with constants independent of the number of
children. The resulting edge-dependent concave recursion creates an
admissible flow by telescoping squared increments, in the spirit of
nonlinear capacity recursions \cite{PP10}.

In the \emph{sufficiency} argument we calibrate deliberately noisy
boundary observations to an arbitrary admissible flow. Unlike the
necessity proof, this requires only elementary identities for the
quadratic information and a signed cubic moment of symmetric
posterior laws. Independent erasure ensures nonnegative information
loss without the imported cone inequality. Sampling a ray from the
flow converts the total loss into a path sum, and the exact
propagation factors leave a kernel bounded by $q^{\text{distance}}$.
This gives a uniform lower bound on root information independently
of the depth and of vertex degrees. Deliberately degraded boundary
observations also appear in robust reconstruction \cite{JM04}; the
present construction calibrates their strengths separately at every
vertex using nonlinear capacity. Because its terminal degradation may be
placed after any fixed additional leaf noise, the same argument proves
robust reconstruction.
